% =====================================================================================================
\section{Closed-form transfer to the four-ball}\label{sec:transfer}

\subsection{Post chart, corner bands, and simplex gluing}\label{sec:transfer-def}

On $\{p<h_K\}$ ($0<h_K<\frac13$, $D=1-3h_K>0$), the post chart $\Pi_{\mathrm{post}}(G,v,w)=(\sigma,\psi,w_{\mathrm{post}})$
with $\psi_0=-\frac16$ and $\ell=1$ sets $u_a=G/P(p)+1$, $\sigma=P(p)v$, $\psi=-\frac16-u_a$, and
$\vartheta=\theta-9(u_a-1)v$ ($w_{\mathrm{post}}=\sqrt{p/\pi}\,e^{2\pi i\vartheta}$). Direct differentiation gives
$d\sigma\wedge d\psi+dp\wedge d\vartheta = (P\,dv-9v\,dp)\wedge(-P^{-1}dG-9GP^{-2}dp)+dp\wedge(d\theta-9\frac{G}{P}dv-9v\frac{dG}{P}-81\frac{Gv}{P^2}dp)
=dG\wedge dv+dp\wedge d\theta$. Under $\Pi_{\mathrm{post}}$, $T_u$ becomes $\tau_u\colon\psi\mapsto\psi-1$, $T_v$
becomes $\tau_v\colon(\sigma,\psi,p,\vartheta)\mapsto(\sigma+P(p),\psi,p,\vartheta+9\psi+\frac12)$ (because $-9\psi_0=\frac32\equiv\frac12\pmod1$),
and the three straight marked fibres $(\mathcal M_1,\mathcal M_2,\mathcal M_3)$ become
\[
 (v_c,\psi)\equiv(0,0),\quad (0,\tfrac12),\quad (\tfrac12,-\tfrac16)\qquad\bigl(v_c=\sigma/P(p),\ \psi\bmod1\bigr).
\]

Fix fractions $(f_m,f_{d_2},f_{\sigma_0},f_\delta)=(0.32,0.20,0.20,0.10)$ of $D$, so $m=0.32D$, $d_2=0.20D$,
$\sigma_0=0.20D$, and $\delta_{\mathrm{kh}}=0.10D$, and let $K_B$ be the corner keyhole on $\bar D(h_K)$ with budget
$0.8\delta_{\mathrm{kh}}$. Set $(\kappa_0,\kappa_1,\kappa_2)=(-\frac1{12},\frac1{12},0)$, phase offsets
$\mathrm{off}_0=(0,0)$, $\mathrm{off}_1=(\frac1{12},\frac14)$, $\mathrm{off}_{2s}=(\frac1{12},\frac14)$,
$\mathrm{off}_{2c}=(\frac1{12},0)$, $\mathrm{left}_0=\mathrm{off}_{2c}$, $\mathrm{left}_1=\mathrm{off}_0$,
$\mathrm{left}_2=\mathrm{off}_1$, and $\kappa=e^{-2\pi i/12}=e^{-i\pi/6}$.

Reduce $(\sigma,\psi,p,\vartheta)$ by $\tau_v^{-n}$ ($n=\lfloor v_c+\frac16\rfloor$) into the fundamental strip
$v_c\in[-\frac16,\frac56)$, and put $c_j=\sigma-j(1-3p)$ for $j\in\{0,1,2\}$. Under the split inequality
$2m<D-3\delta_{\mathrm{kh}}$ ($0.64<0.70$), for any $p_2<h_K+\delta_{\mathrm{kh}}$ we have $1-3p_2>D-3\delta_{\mathrm{kh}}>2m$,
so the three bands $\{|c_j|<m\}$ ($j=0,1,2$) are pairwise disjoint both before and after the corner keyhole!

\begin{enumerate}[label=(\arabic*)]
\item \emph{Corner keyhole $\mathcal B$ (\texttt{Band}).} Outside $\bigcup_{j=0}^2\{|c_j|<m\}$, $\mathcal B=\mathrm{id}$.
On $\{|c_j|<m\}$, put $\tau_B(c)=1-S_4((|c|-d_2)/(m-d_2))$, $k_j(c)=\kappa_j S_4((c+m)/(2m))$, $\psi_1=\psi+a_j+k_j(c)$,
$v_j=\vartheta+b_j+3jk_j(c)-3j\psi_1$ ($(a_j,b_j)=\mathrm{left}_j$), apply $(w_2,g_B)=K_B^{\tau_B(c)}(\sqrt{p/\pi}e^{2\pi iv_j})$,
set $\psi_2=\psi_1-\tau_B'(c)g_B$, $p_2=\pi|w_2|^2$, $v_{j,2}=\arg(w_2)/(2\pi)$, recover
$\sigma_{\mathrm{out}}=c+j(1-3p_2)$ (preserving $c_j$!), $\psi_{\mathrm{out}}=\psi_2-a_j-k_j(c)$,
$\vartheta_{\mathrm{out}}=v_{j,2}+3j\psi_2-b_j-3jk_j(c)$, and restore $\tau_v^n$.
\item \emph{Simplex gluing $\mathcal W$.} On the output of $\mathcal B$ (reduced to $v_c\in[-\frac16,\frac56)$):
 \begin{itemize}
 \item On the four sector pieces outside $\bigcup_{j=0}^2\{|c_j|<m\}$ ($0\le v_c<\frac13$, $\frac13\le v_c<\frac23$,
 $\frac23\le v_c<\frac56$, and $-\frac16\le v_c<0$, indexed by $s\in\{0,1,2s,2c\}$), set $(\tilde\psi,\tilde\vartheta)=(\psi,\vartheta)+\mathrm{off}_s$ and define $(X_0,X_1,X_2;T_1,T_2)$ by
 \begin{align*}
  s=0:\quad &(1-2p-\sigma,\ p+\sigma,\ p;\ \tilde\psi,\ \tilde\vartheta-\tilde\psi),\\
  s=1:\quad &(p,\ 2-5p-\sigma,\ 4p+\sigma-1;\ -\tilde\vartheta+4\tilde\psi,\ -\tilde\vartheta+5\tilde\psi),\\
  s=2s:\quad &(7p+\sigma-2,\ p,\ 3-8p-\sigma;\ \tilde\vartheta-8\tilde\psi,\ -\tilde\psi),\\
  s=2c:\quad &(1-2p+\sigma,\ p,\ p-\sigma;\ \tilde\vartheta+\tilde\psi,\ -\tilde\psi),
 \end{align*}
 and $W=(\sqrt{X_0},\sqrt{X_1}e^{2\pi iT_1},\sqrt{X_2}e^{2\pi iT_2})\in S^5\subset\C^3$.
 \item On band $j\in\{0,1,2\}$ ($|c|<m$, $c=c_j$), put $q(c)=\sigma_0(1-S_4(|c|/d_2))$, $S(c)=\sqrt{c^2+q(c)^2}$,
 $\lambda_+(c)=\frac12(S(c)+c)\ge0$, $\lambda_-(c)=\frac12(S(c)-c)\ge0$ (evaluated without cancellation via
 $q(c)^2/(2(S(c)+|c|))$), $\lambda(c)=\frac12(c-S(c))=-\lambda_-(c)$, $v_\ell=\operatorname{frac}(\vartheta+b_j+3jk_j(c)-3j(\psi+a_j+k_j(c)))\in(0,1)$
 (since $|c|\le d_2\implies\tau_B(c)=1\implies v_\ell\notin\{0,1\}$ on band interiors where $\lambda'\ne0,1$!),
 $x_a=p+\lambda_+(c)$, $x_b=p+\lambda_-(c)$, $x_r=1-x_a-x_b$, $T_a=\psi+a_j+k_j(c)+\lambda'(c)v_\ell$, and permute
 $(\sqrt{x_r},\sqrt{x_a}e^{2\pi iT_a},\sqrt{x_b}e^{2\pi i(v_\ell-T_a)})$ by the cyclic shift $(j,j+1,j+2)\bmod3$.
 \end{itemize}
\item \emph{Rotated ball chart $\chi_\kappa$.} Set $Z_0=(W_2-\kappa W_1)/\sqrt2$, $Z_1=(W_1+\bar\kappa W_2)/\sqrt2$,
$Z_2=W_0$ (unitary rotation in $(W_1,W_2)$, so $|Z|^2=|W|^2=1$). For $Z_0\ne0$, set
$\chi_\kappa([Z])=\frac{\bar Z_0}{\sqrt\pi|Z_0|}(Z_1,Z_2)\in\Bo(1)\subset\C^2\cong\R^4$.
\end{enumerate}

\begin{theorem}[Global regularity, symplecticity, injectivity, and line avoidance of $K_{\mathrm{post}}$]\label{thm:transfer}
Suppose $0<h_K<\frac13$, $K_B$ satisfies \eqref{eq:K123} on $\bar D(h_K)$ with
$\max(\pi d_B^2,\delta_{\mathrm{rel},B},\delta_{\mathrm{small},B})\le\delta_{\mathrm{kh}}$ \emph{(H13)}, $0<d_2<m$,
$2m<D-3\delta_{\mathrm{kh}}$, $\eta_1:=\frac23-(2\delta_{\mathrm{kh}}+S_{\max})/D>0$, and $m\le3\eta_1 D$ \emph{(H14)--(H15)}.
Then on $\mathcal U_{\mathrm{post}}=\{(\sigma,\psi,w_{\mathrm{post}}):p<h_K\}$:
\begin{enumerate}[label=(\roman*)]
\item \emph{Well-definedness and $C^7$ seams.} After $\mathcal B$, $p_2<h_K+\delta_{\mathrm{kh}}$. On every sector piece,
$X_0,X_1,X_2\ge0$ with $X_0+X_1+X_2=1$, and only the coordinate not carrying an angle can vanish (at $p=0$, the angle
matches polar coordinates $w_{\mathrm{post}}=\sqrt{p/\pi}e^{2\pi i\vartheta}$ smoothly). On band $j$, for $|c|\le d_2$
we have $\tau_B(c)=1$, so $p_2>x_{0,B}>0$ and $v_\ell\in(\theta_{0,B},1-\theta_{0,B})$ avoids the branch cut $\{0,1\}$;
for $d_2\le|c|<m$, $q(c)\equiv0$, so $S(c)=|c|$ ($\lambda'(c)=0$ for $c\ge d_2$, $\lambda'(c)=1$ for $c\le-d_2$), making
$e^{2\pi iT_a}$ and $e^{2\pi i(v_\ell-T_a)}$ independent of the integer lift of $v_\ell$ and smooth across $p_2=0$! Across
all six seams $c_j=\pm m$ (where $\tau_B=0$, $k_j(-m)=0$, $k_j(+m)=\kappa_j$) and across the fundamental cut
$v_c=-\frac16\leftrightarrow\frac56$ (via $\mathrm{off}_{2c}-\mathrm{off}_{2s}=(0,-\frac14)$ and $\tau_v$), the
rank-one projector $ZZ^*\in\R^9$ matches identically to order $C^7$ and is invariant under $\tau_u,\tau_v$.
\item \emph{Exact zero set of $Z_0$ (line avoidance).} For $x\in\mathcal U_{\mathrm{post}}$, $Z_0(\mathcal W(\mathcal B(x)))=0$
if and only if $x$ lies in the $(\tau_u,\tau_v)$-orbit of one of the three straight marked fibres
$(0,0)$, $(0,\frac12)$, $(\frac12,-\frac16)$.
\item \emph{Symplectic embedding and cancellation-free inverse.} On $\mathcal U_{\mathrm{post}}\setminus\operatorname{deck}\{\mathcal M_1,\mathcal M_2,\mathcal M_3\}$,
$K_{\mathrm{post}}=\chi_\kappa\circ\mathcal W\circ\mathcal B$ takes values in $\Bo(1)$, satisfies
$K_{\mathrm{post}}^*\omega_0=d\sigma\wedge d\psi+dw_{\mathrm{post},x}\wedge dw_{\mathrm{post},y}$, and identifies two
points if and only if they lie in the same $(\tau_u,\tau_v)$-orbit. Its left inverse (\texttt{transfer.py}) recovers
every sector moment directly from the non-negative coordinate $p=X_2,X_0,X_1$ (never by $1-X_a-X_b$) and every band
moment from $p_2=X_b-\lambda_-(c)$ ($c\ge0$) or $X_a-\lambda_+(c)$ ($c<0$).
\end{enumerate}
\end{theorem}

\begin{proof}
% Clean proof of (i)-(iii) of Theorem 6.1:
(i) On band $j$, $x_r=1-2p_2-S(c)>1-2(h_K+\delta_{\mathrm{kh}})-S_{\max}=\frac13+\frac23D-2\delta_{\mathrm{kh}}-S_{\max}=\frac13+\eta_1 D>\frac13$
by (H14). On each sector, $\min_l|c_l|\ge m>0$ ensures both non-$p$ coordinates are $\ge m>0$ (for instance on sector $1$,
$c_1\ge m\implies X_2=p+c_1\ge m>0$ and $c_2\le-m\implies X_1=-c_2+p\ge m>0$).
At $c_0=-m$: $k_0(-m)=0$, $\mathrm{left}_0=\mathrm{off}_{2c}$, $\lambda(-m)=-m$, and $\lambda'(-m)=1$, giving
$(x_r,x_a,x_b)=(1-2p+m,p,p+m)$ and $(T_a,v_\ell-T_a)=(\tilde\vartheta+\tilde\psi,-\tilde\psi)$, matching sector $2c$;
at $c_0=+m$: $k_0(+m)=\kappa_0=-\frac1{12}$, $\mathrm{left}_0+(\kappa_0,0)=(0,0)=\mathrm{off}_0$, $\lambda(+m)=0$, and
$\lambda'(+m)=0$, giving $(1-2p-m,p+m,p)$ and $(\tilde\psi,\tilde\vartheta-\tilde\psi)$, matching sector $0$.
Similarly $c_1=\pm m$ matches sectors $0\leftrightarrow1$ ($\mathrm{left}_1+(\kappa_1,3\kappa_1)=\mathrm{off}_1$) and
$c_2=\pm m$ matches sectors $1\leftrightarrow2s$ ($\mathrm{left}_2+(0,0)=\mathrm{off}_{2s}$). Under $\tau_v$, a point near
$v_c=-\frac16$ in sector $2c$ maps to $v_c'=v_c+1\in[\frac56,1)$, where substituting $\sigma'=\sigma+3-9p$ and
$\vartheta'=\vartheta+9\psi+\frac12$ into sector $2s$'s formula gives $(X'_0,X'_1,X'_2)=(X_0,X_1,X_2)$, $T'_2=T_2$, and
$T'_1=(\vartheta+9\psi+\frac12+\frac14)-8(\psi+\frac1{12})=\vartheta+\psi+\frac1{12}=T_1\pmod1$.

(ii) Since $|\kappa|=1$, $Z_0=0\iff W_2=\kappa W_1$, which requires $X_1=X_2$ and (when $X_1=X_2>0$)
$T_2-T_1\equiv-\frac1{12}\pmod1$. Across all seven pieces:
\begin{itemize}
\item On band $1$ ($X_1=x_r,X_2=x_a$), $X_1-X_2=1-3p_2-(\frac32S(c)+\frac12c)>0.70D-(\frac32\cdot0.32+0.16)D=+0.06D>0$;
symmetrically $X_2-X_1>+0.06D>0$ on band $2$.
\item On sectors $0$, $2c$, $2s$: $X_1-X_2=\sigma\ge m>0$ ($s=0$), $X_1-X_2=\sigma\le-m<0$ ($s=2c$), and
$X_2-X_1=3-9p-\sigma>\frac16(3-9p)>0$ ($s=2s$).
\item On band $0$ ($X_1=x_a,X_2=x_b$): $X_1=X_2\iff\lambda_+(c)=\lambda_-(c)\iff c=0$ ($v_c=0$). At $c=0$, $\tau_B(0)=1$,
$\tau_B'(0)=0$, $\lambda'(0)=\frac12$, $k_0(0)=-\frac1{24}$, and $\mathrm{left}_0=(\frac1{12},0)$, so
$T_a=\psi+\frac1{24}+\frac12v_\ell$ and $T_2-T_1=v_\ell-2T_a=-2\psi-\frac1{12}$ ($v_\ell$ cancels identically), which equals
$-\frac1{12}\pmod1$ if and only if $\psi\in\{0,\frac12\}\pmod1$ ($\mathcal M_1$ and $\mathcal M_2$).
\item On sector $1$: $X_1=X_2\iff 2-5p-\sigma=4p+\sigma-1\iff v_c=\frac12$ (where $|c_1|=|c_2|=\frac12(1-3p)>\frac12D>m$,
outside all bands). There $X_1=X_2=\frac12(1-p)>\frac13$ and $T_2-T_1=\tilde\psi=\psi+\frac1{12}\equiv-\frac1{12}\pmod1
\iff\psi\equiv-\frac16\pmod1$ ($\mathcal M_3$).
\end{itemize}

(iii) On each sector, $(X_1,X_2)=A_s(p,\sigma)+B_s$ with $A_s\in\mathrm{GL}_2(\Z)$ and
$(T_1,T_2)=(\tilde\vartheta,\tilde\psi)A_s^{-1}$ gives $\sum_{l=1}^2 dX_l\wedge dT_l=dp\wedge d\vartheta+d\sigma\wedge d\psi$;
on band $j$, $dx_a\wedge dT_a+dx_b\wedge d(v_\ell-T_a)=dc\wedge d\psi_2+dp_2\wedge dv_\ell
=d\sigma_{\mathrm{out}}\wedge d\psi_{\mathrm{out}}+dp_2\wedge d\vartheta_{\mathrm{out}}$, and $\mathcal B$ is symplectic
by Proposition~\ref{prop:ks-symp}. Pulling back $\omega_0$ by $\chi_\kappa$ gives $\omega_{\mathrm{FS}}$, whose pullback
by $W$ is $\sum_{l=1}^2dX_l\wedge dT_l$. Finally, the three band image sets $\mathcal R_j=\{|X_{j+1}-X_{j+2}|<m,\ X_j>\frac13\}$
are pairwise disjoint: on $\mathcal R_0$, if $X_1=p_2+\lambda_+(c)>\frac13$, then since $p_2<\frac13-\frac7{30}D$ and
$\max_{|t|\le d_2}\lambda_+(t)=0.20D<\frac7{30}D$, we must have $c>d_2$, so $S(c)=c$, $\lambda_-(c)=0$, and
$X_0-X_2=1-3p_2-c>\frac7{10}D-0.32D=0.38D>m$; symmetrically $X_0-X_1>0.38D>m$ whenever $X_2>\frac13$. Outside
$\mathcal R_0\cup\mathcal R_1\cup\mathcal R_2$, $\operatorname{argmin}(X_0,X_1,X_2)$ uniquely identifies sector $0$
($X_2$ minimal), sector $1$ ($X_0$ minimal), or sector $2s$/$2c$ ($X_1=p$ minimal: on $2c$, $X_2\ge p+m>p$ and
$X_0-X_1=1-3p+\sigma\ge\frac12(1-3p)>0$; since $2c$ is the $\tau_v^{-1}$-image of $v_c'\in[\frac56,1)$ where $2s$'s
formula produces the identical $W$, $2s$'s inverse followed by \texttt{reduce} recovers $(v_c,\psi,p,\vartheta)$ exactly).
\end{proof}

% =====================================================================================================
\section{Explicit parameter verification and proof of Theorem~\ref{thm:main}}\label{sec:assembly}

\subsection{The closed-form parameter tuple $\Pi^*(k,\eps_c;\frac14)$}\label{sec:pistar}

\begin{definition}[Parameter tuple $\Pi^*(k,\eps_c;f)$]\label{def:pistar}
For an integer $k\ge10$, $\eps_c\in(0,\frac1{100}]$, and $f=\frac14$, put $s=\sqrt k$,
$\eta=\eps_c/(1+\sqrt{1-\eps_c})$, $t_0=1-\frac12\eta$, $a=t_0/s$, $h=t_0/3$, $c_s=3s-9$, $g=\frac32\eta$, $T=g/c_s$,
$(\lambda_a,\lambda_2,\mu)=(\frac9{10},\frac3{10},\frac12)$, $w_G=\lambda_a g$, $w_2=\lambda_2 g$, $\upsilon=\sqrt{11/1000}$.
\begin{enumerate}[label=(\arabic*)]
\item \emph{Comparator keyhole $K_C$:} $r_a=\sqrt{a/\pi}$, $d=\sqrt{fT/\pi}/(1+\upsilon)$,
$\beta=fT/\bigl(\pi(r_a+d)(r_a+8.4d)\bigr)$, $W_1=\beta/(3(1+\beta))$, $r_h=0.9d\sin(2\pi W_1)$,
$\rho_1=0.8((d-r_h)/(d+r_a))^2$, $W_2=\beta/(1+\beta-\rho_1)-W_1$, $X_w=0.2\pi(d-r_h)^2$.
\item \emph{Post height $h_K$ and $D=1-3h_K$:}
$h_K=\max\bigl(a,\ O(-w_G)-\frac12\mu T(1-\lambda_a),\ h-\frac12\mu T\bigr)$, so
$D=\min\bigl(\frac{s-3+\frac32\eta}{s},\ \frac{s-3+w_G}{s}+\frac32\mu T(1-\lambda_a),\ \frac12\eta+\frac32\mu T\bigr)>0$.
\item \emph{Corner keyhole $K_B$:} $(m,d_2,\sigma_0,\delta_{\mathrm{kh}})=(0.32,0.20,0.20,0.10)D$,
$r_B=\sqrt{h_K/\pi}$, $d_B=\sqrt{0.8\delta_{\mathrm{kh}}/\pi}/1.05$,
$\beta_B=0.8\delta_{\mathrm{kh}}/\bigl(\pi(r_B+d_B)(r_B+8.4d_B)\bigr)$, with $(W_{1,B},r_{h,B},\rho_{1,B},W_{2,B},X_{w,B})$
by the same `keyhole\_star` formulas with $\rho_1$-factor $0.7$.
\end{enumerate}
\end{definition}

\begin{lemma}[Grid-free analytic bounds for (H6) and (H7)]\label{lem:h67}
With $(\lambda_a,\lambda_2,\mu,f)=(\frac9{10},\frac3{10},\frac12,\frac14)$, suppose $\delta_{\mathrm{rel}}=fT=\frac14T$
and $\delta_{\mathrm{small}}\le0.04T$. Then (H6) and (H7) hold on all of $[0,1]$ with margin $>0.012T$.
\end{lemma}

\begin{proof}
Divide by $T$. For (H6), $(1-\mu)(1-\lambda_a+u(\lambda_a-\lambda_2))=0.05+0.30u\ge0.05>0.04\ge\delta_{\mathrm{small}}/T$.
It remains to bound $h_6(u)=0.05+0.30u-0.25S_4(u)$ on $[0,1]$. On $[0,\frac12]$, Lemma~\ref{lem:steps}(i) gives
$S_4(u)\le u$, so $h_6(u)\ge0.05+0.05u\ge0.05$. On $[\frac12,1]$, $S_4$ is concave, so for $u_0=\frac23$
($S_4(\frac23)<0.9414$, $S_4'(\frac23)=\frac{256}{70}\sin^8(\frac{2\pi}{3})=\frac{81}{70}<1.1572$),
$S_4(u)\le S_4(u_0)+S_4'(u_0)(u-u_0)$, bounding $h_6(u)$ from below by an affine function whose values at $u=\frac12$
and $u=1$ are $>0.0128$ and $>0.0182$.
For (H7), since $\max(\delta_{\mathrm{rel}},\delta_{\mathrm{small}})/T=\frac14$, we must bound
$h_7(u)=0.5(0.7+0.3u)-0.25-0.3I_4(u)=0.10+0.15u-0.30I_4(u)$ on $[0,1]$. Its derivative $0.30(\frac12-S_4(u))$ is
$\ge0$ on $[0,\frac12]$ and $\le0$ on $[\frac12,1]$, so $\min_{[0,1]}h_7=\min(h_7(0),h_7(1))=0.10>0$.
\end{proof}

\begin{proposition}[Universal validity of (H0)--(H15)]\label{prop:all-hyps}
For every integer $k\ge10$ and every $\eps_c\in(0,\frac1{100}]$, $\Pi^*(k,\eps_c;\frac14)$ satisfies all hypotheses of
Theorems~\ref{thm:kh-bounds}, \ref{thm:comparator}, \ref{prop:twists}, and \ref{thm:transfer}, as well as $r_+>1$.
\end{proposition}

\begin{proof}
Since $s\ge\sqrt{10}>3.16227$ and $0<\eta\le1-\sqrt{0.99}<0.005013$, we have $t_0>0.9974$,
$s/(s-3)\le\sqrt{10}/(\sqrt{10}-3)=10+3\sqrt{10}<19.49$, and
\[
 \frac{T}{a} = \frac{\frac12\eta/(s-3)}{t_0/s} = \frac{\eta}{2t_0}\,\frac{s}{s-3}
 < \frac{0.005013}{1.9948}\cdot 19.49 < 0.049.
\]
\emph{(H2) and $r_+>1$:} Since $0<m_s\le\frac18\eta$ and $\eta<0.01$, $(1-m_s)^2>1-2m_s\ge1-\frac14\eta$, while
$(1-\frac12\eta)/(1-\eta)=1+\frac{\eta}{2(1-\eta)}>1+\frac12\eta$. Thus
$r_+^2=(1-m_s)^2(1-\frac12\eta)/(1-\eta)>(1-\frac14\eta)(1+\frac12\eta)=1+\frac14\eta(1-\frac12\eta)>1$, which is
equivalent to $\rho_b=\rho_{\mathrm{ball}}/(1-m_s)^2<a$.

\emph{(H3)--(H5) for $K_C$:} Since $\pi(r_a+d)(r_a+8.4d)>\pi r_a^2=a$, $\beta<\frac14(T/a)<0.0123$. Also
$\pi d^2=\frac14T/(1+\upsilon)^2<0.205T$, so $(d/r_a)^2<0.205(0.049)<0.0101$ ($d/r_a<0.101$). By definition, (K1) and
(K2) hold with equality/$0.9<1$, $r_h<0.9d(2\pi W_1)<0.024d<d$, $\rho_1<0.8(0.101/1.101)^2<0.0068<1$,
$W_2-W_1=\beta\bigl(\frac1{1+\beta-\rho_1}-\frac2{3(1+\beta)}\bigr)\ge\frac{\beta}{3(1+\beta)}>0$, and
$W_2<W_1+W_2=\beta/(1+\beta-\rho_1)<0.0123/0.993<0.013<\frac16$. For (K3),
$\rho_1\pi(d+r_a)^2+\frac12X_w=(0.8+0.1)\pi(d-r_h)^2=0.9\pi(d-r_h)^2<\pi(d-r_h)^2$ holds identically! Moreover,
$d/(r_a+8.4d)<\frac1{8.4}$ and $W_2/\beta<\frac1{0.993}-\frac1{3.04}<0.68$, so
\[
 \frac{\delta_{\mathrm{small}}}{fT} = 4\pi^2\,\frac{d}{r_a+8.4d}\,\Bigl(\frac{W_2}{\beta}\Bigr)^2\beta
 < \frac{4\pi^2}{8.4}\cdot(0.68)^2\cdot 0.0123 < 0.027 < 0.04.
\]
Lemma~\ref{lem:h67} then gives (H6) and (H7).

\emph{(H8), (H9), (H9b):} $\pi d^2<0.205T<0.35T=T(1-\mu-\frac12\lambda_2)$, and
$s(\pi d^2+0.65T)<0.855sT=0.4275\eta\frac{s}{s-3}<0.042<1-0.45\eta=sO(-w_G)$, giving (H8). For (H9),
$c_s\pi r_h^2<c_s\pi(0.024d)^2<6\cdot10^{-4}c_s(0.205T)<2\cdot10^{-4}g$, so $c_s\pi r_h^2+1.5\lambda_2g<(0.45+0.001)g<g$.
For (H9b), $a<1/\sqrt{10}\implies 3-9a>3(1-9/\sqrt{10})>0.1539$, while $\delta=\frac14w_2=\frac{9}{80}\eta<0.0006<0.1539$.
(H10) is a check of fixed rational constants ($0.045+0.045<0.10$, $0.045+0.14<0.25$, $2\cdot0.12<\frac14$,
$\frac16-0.12-0.02>0.02>0$).

\emph{(H11)--(H15):} By Definition~\ref{def:pistar}(2), $D>0$ ($h_K<\frac13$), $h_K\ge a$, and since $O(G)-\mu\operatorname{slack}(G)$
is piecewise affine on $[-w_G,0]\cup[0,G_{\max}]$, its maximum is bounded by $\max\bigl(O(-w_G)-\mu T(1-\lambda_a),O(0)-\mu T,O(G_{\max})-\mu T\bigr)$;
note $O(0)<O(G_{\max})=h$ when $G_{\max}>0$ ($O(-w_G)$ otherwise), and subtracting $\frac12\mu T(1-\lambda_a)>0$ in $h_K$
ensures strict inequality $h_K>O(G)-\mu\operatorname{slack}(G)$ everywhere on $[-w_G,G_{\max}]$.
Finally, $h_K\ge a>\frac{0.997}{s}$ while $\delta_{\mathrm{kh}}=0.1D$: since $D=1-3h_K\le 1-3a<1$ and even as $s\to\infty$,
$h_K\ge\max(a,\frac13-\frac13(\frac12\eta+\frac32\mu T))>\frac13-0.01>0.32$, we always have $h_K>0.31$ (for $k=10$,
$a>0.315$; for $k\ge11$, $O(-w_G)$ or $h=\frac13t_0>0.332$ gives $h_K>0.315$!). Thus $\delta_{\mathrm{kh}}/h_K<0.10/0.315<0.32$
and $\beta_B<0.8(\delta_{\mathrm{kh}}/h_K)/(1)<0.003$ (since $D\le\frac12\eta+\frac32\mu T<0.015\implies\delta_{\mathrm{kh}}<0.0015$).
Exactly as for $K_C$, $K_B$ satisfies (K1)--(K3) with $\rho_{1,B}\pi(d_B+r_B)^2+\frac12X_{w,B}=0.8\pi(d_B-r_{h,B})^2$,
$\delta_{\mathrm{rel},B}=0.8\delta_{\mathrm{kh}}$, $\pi d_B^2=0.8\delta_{\mathrm{kh}}/1.05^2<\delta_{\mathrm{kh}}$, and
$\delta_{\mathrm{small},B}<0.01\delta_{\mathrm{kh}}$, proving (H12)--(H13). On $[-m,m]$, for $|c|\le d_2$ we have
$S(c)\le\sqrt{d_2^2+\sigma_0^2}=0.20\sqrt2\,D<0.283D<0.32D=m$, and for $d_2\le|c|\le m$, $S(c)=|c|\le m$; thus
$S_{\max}=0.32D$, $\eta_1=\frac23-(0.20+0.32)=\frac{11}{75}>0.1466>0$, $3\eta_1=\frac{11}{25}=0.44\ge0.32=m/D$ (H14--H15),
and $2m/D+3\delta_{\mathrm{kh}}/D=0.64+0.30=0.94<1$ (`split`).
\end{proof}

\subsection{Completion of the proof of Theorem~\ref{thm:main}}\label{sec:proof-end}

\begin{proof}[Proof of Theorem~\ref{thm:main}]
Fix $k\ge10$ and $\eps\in(0,1)$. Since $\eps_c=\min(\eps,\frac1{100})\le\eps$, $R_c\le R_{\mathrm{req}}$, and
Proposition~\ref{prop:all-hyps} gives $r_+>1$ and all hypotheses (H0)--(H15).
\begin{enumerate}[label=(\arabic*)]
\item Let $X=(q_1,q_2,p_1,p_2)\in\EB^{\circ}(r_+)$. Stage~1, $S_{\mathrm{in}}(X)=(q_1,p_1,q_2,p_2)/(\sqrt\pi R_s)$,
satisfies $\pi|S_{\mathrm{in}}(X)|^2=|X|^2/R_s^2<r_+^2\rho_b=a$ and pulls back $dx\wedge dy+\omega_V$ to
$(\pi R_s^2)^{-1}\omega_0$.
\item By Proposition~\ref{prop:kr}, $\mathrm{KR}$ maps $\Bo(a)$ real-analytically, bijectively, and symplectically onto
$\mathcal C_a$.
\item By Lemma~\ref{lem:cell-flat} and Theorem~\ref{thm:comparator}, $F^{(i)}\circ C_i$ ($0\le i<k$) are $C^7$
symplectic embeddings on $\mathcal C_a$ whose images $U_i=F^{(i)}(C_i(\mathcal C_a))$ are pairwise disjoint modulo
$\langle T_u,T_v\rangle$, lie in $\{p<h_K\}$, and miss the corridor $\mathcal K_{\mathrm{free}}$.
\item By Proposition~\ref{prop:twists}, $\Psi=\Phi_g\circ\Psi_S$ is a $C^4$ deck-equivariant symplectomorphism preserving
$\{p<h_K\}$ and mapping $\Psi(U_i)$ into $\{p<h_K\}\setminus\operatorname{deck}\{\mathcal M_1,\mathcal M_2,\mathcal M_3\}$,
keeping $\Psi(U_0),\dots,\Psi(U_{k-1})$ pairwise disjoint modulo $\langle T_u,T_v\rangle$.
\item By Theorem~\ref{thm:transfer}, $K_{\mathrm{post}}\circ\Pi_{\mathrm{post}}$ is $C^7$, symplectic, deck-invariant,
and injective modulo $\langle T_u,T_v\rangle$ on $\{p<h_K\}\setminus\operatorname{deck}\{\mathcal M_1,\mathcal M_2,\mathcal M_3\}$
into $\Bo(1)$. Hence $\tilde f_i=K_{\mathrm{post}}\circ\Pi_{\mathrm{post}}\circ\Psi\circ F^{(i)}\circ C_i\circ\mathrm{KR}$
are $C^4$ symplectic embeddings on $\Bo(a)$ with pairwise disjoint images in $\Bo(1)$.
\item Finally, $S_{\mathrm{out}}(y_1,y_2,y_3,y_4)=\sqrt\pi R_s(y_1,y_2,y_3,y_4)$ maps $\Bo(1)$ onto $\EB^{\circ}(R_s)\subset\EB^{\circ}(R_{\mathrm{req}})$
with conformal factor $\pi R_s^2$, cancelling Stage~1's $(\pi R_s^2)^{-1}$ so that $f_i^*\omega_0=\omega_0$.
\item For $Y\in\R^4$, $\mathrm{Preimage}(j,Y)$ checks $|Y|<R_s$ (via $1-\pi|Y/(\sqrt\pi R_s)|^2>0$ in `unchart`), applies
$\mathrm{to\_fund}\circ\Psi^{-1}\circ\Pi_{\mathrm{post}}^{-1}\circ\mathcal B^{-1}\circ(\chi_\kappa\circ\mathcal W)^{-1}$
to obtain the unique fundamental representative $y_{\mathrm{fl}}\in\mathcal F_{\mathrm{fl}}$, applies
$(F^{(j)})^{-1}$ and $C_j^{-1}$ to obtain $(X_{\mathrm{cell}},T_{\mathrm{cell}},w)$, tests membership in $\mathcal C_a$
($p<a$, $0<X_{\mathrm{cell}}<a-p$, $0<T_{\mathrm{cell}}<1$), and applies $S_{\mathrm{in}}^{-1}\circ\mathrm{KR}^{-1}$.
When $Y=f_i(X)$ with $X\in\EB^{\circ}(r_+)$, deck invariance of $K_{\mathrm{post}}\circ\Pi_{\mathrm{post}}$ and
Theorem~\ref{thm:comparator}(ii) guarantee that $\mathrm{to\_fund}$ recovers the exact point $(F^{(i)}\circ C_i\circ\mathrm{KR}\circ S_{\mathrm{in}})(X)\in\mathcal F_{\mathrm{fl}}$,
giving $\mathrm{Preimage}(i,f_i(X))=(X,\mathrm{true})$ and disjointness for $j\ne i$.\qedhere
\end{enumerate}
\end{proof}

\begin{remark}[Exact arithmetic versus IEEE-754 \texttt{float64}]\label{rem:float64}
Theorem~\ref{thm:main} holds over $\R$ for all $k\ge10$ and $\eps\in(0,1)$. In \texttt{float64}, \texttt{solution(k, eps)}
evaluates the analytic margins of Proposition~\ref{prop:all-hyps} on the rounded parameters and raises \texttt{ValueError}
if machine precision underflows a strict inequality (for instance $\eps\le10^{-16}$, where \texttt{1.0 - eps == 1.0}
makes $a-\rho_b=0$). Within the \texttt{float64}-admissible range, all three scalar Newton solvers run $60$ safeguarded
iterations on bounded brackets (reaching $53$-bit accuracy in $\le8$ steps once quadratic), and the inverse transfer
recovers every moment coordinate without subtraction cancellation. Because gaps of size $O(\eps_c/\sqrt k)$ sit next to
$O(1)$ coordinates, local Lipschitz constants reach $10^8$--$10^{15}$ for small $\eps$ or large $k$; arbitrary-precision
evaluation via \texttt{program/check\_embedding\_mp.py} verifies the submitted Python source directly to $100+$ digits.
\end{remark}

\section*{Authorship and the role of AI}

Liang Song is the sole author and is responsible for the manuscript and its release. The mathematics, the code and the
text were produced by AI systems, including GPT-6.1 Sol (OpenAI), Claude Opus 5.5 (Anthropic), Gemini (Google) and
other models, working in a research environment designed and directed by Liang Song, who chose the problem, assembled
the context, designed the verification pipeline and made the editorial decisions. The theorems of this paper are not
part of the Lean~4 formalisation of Papers~I--II; numerical and arbitrary-precision checks are documented in
\texttt{program/README.md}.
